\section*{ID7973: Classification: (E, I, σ)}
\paragraph{Metadata.}
PiID: 5626476691490 | RTEID: RTE:P30960.9296:T0.7486 | UUID: b707ae26-6f69-5166-afd7-b8407f700c65

\paragraph{Canonical Statement.}
The next step is not “more simulation,” it’s **classification**: which resonances survive in which regions of \textbackslash{}((e, I, \textbackslash{}sigma)\textbackslash{})-space, and why. Once that table exists, comparison to exoplanet data becomes almost boring — which is exactly what you want when the physics is right.

\paragraph{Canonical Equation.}
\[
(e, I, \sigma)
\]

\paragraph{Dictionary.}
The next step is not “more simulation,” it’s classification : which resonances survive in which regions of \textbackslash{}((e, I, )\textbackslash{})-space, and why.

\paragraph{Encyclopedia.}
Classification: (E, I, σ) is a symbolic expression, written (e, I, σ). It introduces a symbol or relation used elsewhere in the framework. It is built from σ, I. Within the Trawin Zero Point registry it serves as one element of the framework's mathematical narrative.
